\documentclass[11pt]{article}
\usepackage[margin=1in]{geometry}
\usepackage{amsmath,amssymb,amsthm}
\usepackage{physics}
\usepackage{enumitem}
\usepackage{hyperref}
\usepackage{titlesec}

\theoremstyle{definition}
\newtheorem{definition}{Definition}[section]
\newtheorem{example}{Example}[section]
\newtheorem{remark}{Remark}[section]
\newtheorem{proposition}{Proposition}[section]
\newtheorem{theorem}{Theorem}[section]

\titleformat{\section}{\large\bfseries}{Lecture 4.\thesection}{1em}{}

\title{\textbf{Lecture 4: Entanglement in Action} \\ \large Understanding Quantum Information \& Computation}
\author{Based on the lectures by John Watrous (IBM Quantum)}
\date{}

\begin{document}
\maketitle
\tableofcontents
\newpage

%----------------------------------------------------------------------
\section{Overview}
%----------------------------------------------------------------------

This lecture concludes the first unit (Basics of Quantum Information) by examining three foundational examples in which \textbf{entanglement} plays a central role:

\begin{enumerate}
    \item \textbf{Quantum Teleportation} --- transmitting a qubit using classical communication and shared entanglement.
    \item \textbf{Superdense Coding} --- transmitting two classical bits via one qubit and shared entanglement.
    \item \textbf{The CHSH Game} --- a non-local game demonstrating that entanglement enables correlations impossible in any classical theory.
\end{enumerate}

These are not merely illustrative examples; they are cornerstones of quantum information theory.

%----------------------------------------------------------------------
\section{Preliminaries: Alice, Bob, and Entanglement}
%----------------------------------------------------------------------

\paragraph{Alice and Bob.}
Hypothetical agents (sender/receiver, cooperating players, etc.) conventionally labelled $A$ and $B$.  They may be in different locations; the setup specifies what resources they share and what communication is permitted.

\paragraph{The e-bit.}
When we say Alice and Bob \emph{share an e-bit}, we mean Alice holds qubit $A$, Bob holds qubit $B$, and the joint state is
\[
    \ket{\Phi^+}_{AB} = \frac{1}{\sqrt{2}}\bigl(\ket{00}+\ket{11}\bigr).
\]
This represents \textbf{one unit of entanglement}.  Compare with \emph{one bit of shared randomness}: a pair of classically correlated bits in state $(1/2,0,0,1/2)^T$, which is \emph{not} entangled.

\paragraph{Classical correlation vs.\ entanglement.}
Intuitive descriptions (``measuring one instantly affects the other'') do not distinguish $\ket{\Phi^+}$ from classical correlation.  The three examples below \emph{do} distinguish them: each achieves something provably impossible with classical correlations alone.

%----------------------------------------------------------------------
\section{Quantum Teleportation}
%----------------------------------------------------------------------

\subsection{The Scenario}

\begin{itemize}
    \item Alice holds a qubit $Q$ in an unknown state $\ket{\psi}=\alpha\ket{0}+\beta\ket{1}$.
    \item She wishes to transmit $\ket{\psi}$ to Bob.
    \item \textbf{Constraint:} Alice can only send \emph{classical} bits to Bob.
    \item \textbf{Resource:} Alice and Bob share one e-bit ($A$ with Alice, $B$ with Bob).
\end{itemize}

Without the e-bit, transmitting a qubit via classical communication alone is \emph{impossible} (a consequence related to the no-cloning theorem).

\subsection{The Protocol}

\begin{enumerate}
    \item Alice applies CNOT with $Q$ as control, $A$ as target.
    \item Alice applies $H$ to $Q$.
    \item Alice measures both $Q$ and $A$ in the standard basis, obtaining bits $b$ and $a$ respectively.
    \item Alice sends $(a,b)$ to Bob (2 classical bits).
    \item Bob applies $X^a Z^b$ to his qubit $B$.
\end{enumerate}

Bob's correction operations:
\[
\begin{array}{c|c}
(a,b) & \text{Bob's operation} \\\hline
(0,0) & I \\
(0,1) & Z \\
(1,0) & X \\
(1,1) & XZ
\end{array}
\]

\subsection{Analysis}

The initial three-qubit state (ordering: $B, A, Q$) is
\[
    \ket{\Phi^+}_{BA}\otimes\ket{\psi}_Q
    = \frac{1}{\sqrt{2}}\bigl(\ket{00}+\ket{11}\bigr)_{BA}\otimes\bigl(\alpha\ket{0}+\beta\ket{1}\bigr)_Q.
\]

\paragraph{After CNOT ($Q\to A$):}
\[
    \frac{1}{\sqrt{2}}\bigl(\alpha\ket{000}+\alpha\ket{110}+\beta\ket{011}+\beta\ket{101}\bigr).
\]

\paragraph{After Hadamard on $Q$:}
Expand and regroup by isolating the standard basis states of $(A,Q)$:
\[
    \frac{1}{2}\Bigl[
        \ket{00}_{AQ}\otimes(\alpha\ket{0}+\beta\ket{1})_B
        +\ket{01}_{AQ}\otimes(\alpha\ket{0}-\beta\ket{1})_B
        +\ket{10}_{AQ}\otimes(\alpha\ket{1}+\beta\ket{0})_B
        +\ket{11}_{AQ}\otimes(\alpha\ket{1}-\beta\ket{0})_B
    \Bigr].
\]

\paragraph{Key observations:}
\begin{enumerate}
    \item Each measurement outcome $(a,b)$ occurs with probability $\tfrac{1}{4}$, \emph{independent} of $\alpha,\beta$.
    \item Bob's qubit after outcome $(a,b)$ is $X^a Z^b\ket{\psi}$ (up to normalisation).
    \item Applying $X^a Z^b$ corrects to $\ket{\psi}$.
\end{enumerate}

\subsection{Teleportation with Entangled Inputs}

If $Q$ is entangled with an external system $R$, a similar (but slightly more involved) calculation shows the protocol still works perfectly.  By linearity, once correctness is established for arbitrary $\ket{\psi}=\alpha\ket{0}+\beta\ket{1}$, it extends to entangled inputs automatically.

\subsection{Remarks}

\begin{itemize}
    \item Teleportation is a \emph{protocol}, not an ``application.''  It implements quantum communication from classical communication + entanglement.
    \item The e-bit is consumed (``burned up'') during the protocol.
    \item Cost: 1 e-bit + 2 classical bits $\longrightarrow$ 1 qubit of quantum communication.
    \item No information about $\ket{\psi}$ is revealed by the measurement outcomes (uniform distribution).
    \item Alice no longer holds $\ket{\psi}$ after the protocol, consistent with no-cloning.
\end{itemize}

%----------------------------------------------------------------------
\section{Superdense Coding}
%----------------------------------------------------------------------

\subsection{The Scenario}

\begin{itemize}
    \item Alice has two classical bits $(a,b)$ she wishes to send to Bob.
    \item \textbf{Resource:} One e-bit (qubit $A$ with Alice, qubit $B$ with Bob).
    \item Alice may send \emph{one qubit} to Bob (but no classical bits).
\end{itemize}

By \textbf{Holevo's theorem}, transmitting two classical bits via one qubit is impossible \emph{without} entanglement.

\subsection{The Protocol}

\begin{enumerate}
    \item Depending on $(a,b)$, Alice applies an operation to her qubit $A$:
    \[
    \begin{array}{c|c|c}
    (a,b) & \text{Alice's gate} & \text{Resulting Bell state} \\\hline
    (0,0) & I   & \ket{\Phi^+} \\
    (0,1) & Z   & \ket{\Phi^-} \\
    (1,0) & X   & \ket{\Psi^+} \\
    (1,1) & ZX  & \ket{\Psi^-}
    \end{array}
    \]
    \item Alice sends qubit $A$ to Bob.
    \item Bob applies CNOT ($A$ control, $B$ target), then $H$ on $A$, then measures both qubits.  The measurement outcome is $(a,b)$.
\end{enumerate}

\subsection{Analysis}

The key insight is that Alice, by acting on her qubit alone, can steer the shared Bell state to \emph{any} of the four Bell states.  Since the Bell states form an orthonormal basis, Bob can perfectly distinguish them.

Bob's decoding circuit (CNOT then $H$) is the inverse of the Bell-state preparation circuit from Lecture~3.  It maps Bell states back to standard basis states:
\[
    \ket{\Phi^+}\mapsto\ket{00},\quad
    \ket{\Phi^-}\mapsto\ket{01},\quad
    \ket{\Psi^+}\mapsto\ket{10},\quad
    \ket{\Psi^-}\mapsto\ket{11}.
\]

\subsection{Remarks}

\begin{itemize}
    \item Cost: 1 e-bit + 1 qubit of quantum communication $\longrightarrow$ 2 classical bits.
    \item Teleportation and superdense coding together establish an equivalence:
    \[
        1\text{ qubit} + 1\text{ e-bit} \;\longleftrightarrow\; 2\text{ classical bits} + 1\text{ e-bit}.
    \]
    \item Without the e-bit, neither conversion is possible.
\end{itemize}

%----------------------------------------------------------------------
\section{The CHSH Game}
%----------------------------------------------------------------------

\subsection{Non-Local Games}

A \textbf{non-local game} involves:
\begin{itemize}
    \item Two cooperating players, Alice and Bob, who \emph{cannot communicate} during the game.
    \item A \textbf{referee} who sends questions and evaluates answers.
\end{itemize}

Procedure:
\begin{enumerate}
    \item Referee sends question $x$ to Alice and $y$ to Bob (chosen randomly).
    \item Alice responds with $a$; Bob responds with $b$.
    \item Referee declares win/loss based on a fixed, publicly known rule applied to $(x,y,a,b)$.
\end{enumerate}

Alice and Bob may agree on any strategy beforehand (including sharing entanglement).

\subsection{CHSH Game Specification}

\begin{itemize}
    \item Questions and answers are all single bits.
    \item $(x,y)$ is chosen uniformly at random from $\{0,1\}^2$.
    \item \textbf{Winning condition:} $a\oplus b = x\wedge y$, i.e.:
\end{itemize}
\[
\begin{array}{c|c}
(x,y) & \text{Win iff} \\\hline
(0,0) & a = b \\
(0,1) & a = b \\
(1,0) & a = b \\
(1,1) & a \neq b
\end{array}
\]

\subsection{Classical Strategies}

\begin{theorem}
No classical (deterministic or probabilistic) strategy can win the CHSH game with probability exceeding $\frac{3}{4}$.
\end{theorem}

\begin{proof}[Proof sketch]
A deterministic strategy defines functions $a(x)$ and $b(y)$.  The four winning conditions yield four equations in four binary unknowns.  Taking the XOR of all four left-hand sides gives $0$, but the XOR of all four right-hand sides gives $1$---contradiction.  Hence at least one equation must fail, so at most 3 out of 4 cases can be won.

Probabilistic strategies are convex combinations of deterministic ones, so the winning probability is at most the maximum over deterministic strategies, which is $\tfrac{3}{4}$.
\end{proof}

\subsection{Quantum Strategy}

Define the parameterised state and unitary:
\[
    \ket{\chi_\theta} = \cos\theta\,\ket{0}+\sin\theta\,\ket{1},
    \qquad
    U_\theta = \begin{pmatrix}\cos\theta & -\sin\theta \\ \sin\theta & \cos\theta\end{pmatrix}.
\]

\paragraph{Strategy.}
Alice and Bob share $\ket{\Phi^+}$.
\begin{itemize}
    \item Alice applies $U_0=I$ if $x=0$, or $U_{\pi/4}$ if $x=1$, then measures.
    \item Bob applies $U_{\pi/8}$ if $y=0$, or $U_{-\pi/8}$ if $y=1$, then measures.
\end{itemize}

\paragraph{Analysis.}
Using the identity
\[
    \bigl(\ket{\chi_\alpha}\otimes\ket{\chi_\beta}\bigr)^\dagger\ket{\Phi^+}
    = \frac{1}{\sqrt{2}}\cos(\alpha-\beta),
\]
one computes the probability that the measurement outcomes agree or disagree.  In every case, the winning probability is
\[
    p_{\text{win}} = \frac{2+\sqrt{2}}{4} \approx 85.36\%.
\]

\begin{theorem}[Tsirelson's inequality]
No quantum strategy can win the CHSH game with probability exceeding $\frac{2+\sqrt{2}}{4}$.  The strategy above is therefore \textbf{optimal}.
\end{theorem}

\subsection{Geometric Interpretation}

The four pairs of angles $(\alpha,\beta)$ are arranged so that in the first three cases ($x\wedge y=0$), the angle difference $|\alpha-\beta|=\pi/8$ is small (high probability of equal outcomes), while in the fourth case ($x\wedge y=1$), the angle difference is $3\pi/8$ (high probability of unequal outcomes).  This is precisely the structure needed to maximise the overall winning probability.

\subsection{Significance}

\begin{itemize}
    \item The CHSH game provides a concrete, \emph{testable} demonstration of Bell's theorem: quantum mechanics is incompatible with local hidden-variable theories.
    \item Experimental Bell tests (Aspect, Clauser, Zeilinger --- 2022 Nobel Prize in Physics) confirm that quantum strategies achieve the predicted advantage.
    \item The CHSH game is an instance of the broader and actively researched theory of \textbf{non-local games}.
\end{itemize}

%----------------------------------------------------------------------
\section{Summary}
%----------------------------------------------------------------------

\begin{center}
\begin{tabular}{lccc}
\hline
\textbf{Protocol/Game} & \textbf{Input} & \textbf{Output} & \textbf{Cost} \\\hline
Teleportation & 1 qubit & 1 qubit (at Bob) & 1 e-bit + 2 cbits \\
Superdense coding & 2 cbits & 2 cbits (at Bob) & 1 e-bit + 1 qubit \\
CHSH game & --- & win probability & 1 e-bit; $p_{\max}^Q=\frac{2+\sqrt{2}}{4}>\frac{3}{4}=p_{\max}^C$ \\
\hline
\end{tabular}
\end{center}

All three examples demonstrate capabilities of entanglement that are provably impossible with classical resources alone, establishing entanglement as a genuine physical resource.

\end{document}
